\documentclass[12pt]{article}

% Load your custom style package
\usepackage{style}
\usepackage{macros}
\usepackage{amsmath}
\usepackage{subfiles}

%\usepackage{graphicx}
%\usepackage{pgf,tikz,pgfplots}
%\usepackage{tikz-cd}
%\usepackage{mathrsfs}
%\usepackage{mathpple}
%\usepackage{CJKutf8}
%\usepackage{amsmath}
%\usepackage{tikz}
%\usepackage{mathdots}
%%\usepackage{yhmath}
%\usepackage{cancel}
%%\usepackage{color}
%\usepackage{siunitx}
%\usepackage{array}
%\usepackage{multirow}
%\usepackage{amssymb}
%\usepackage{gensymb}
%\usepackage{tabularx}
%\usepackage{booktabs}
%\usetikzlibrary{fadings}
%\usetikzlibrary{patterns}
%\usetikzlibrary{shadows.blur}
%\usetikzlibrary{shapes}

\usepackage[a4paper,margin=1.2in]{geometry}
\usepackage{fancyhdr}
\pagestyle{fancy}
\fancyhf{}
%\fancyhead[L]{\headfam\headweight Quiz}
%\fancyhead[C]{\headfam\itshape\large\thetitle}
\fancyhead[R]{\headfam\headweight page \thepage}

%bib source
\addbibresource{RRChiral.bib}


% Title formatting uses your package's redefinition, so normal \title works
\title{Higher-dimensional Virasoro algebras}
\author{Zhengping Gui and Brian R. Williams}
\date{\today}


% MATH for zg -----------------------------------------------------------
\renewcommand{\norm}[1]{\left\Vert#1\right\Vert}
\renewcommand{\abs}[1]{\left\vert#1\right\vert}
\newcommand{\set}[1]{\left\{#1\right\}}
\newcommand{\Real}{\mathbb R}
%\newcommand{\eps}{\varepsilon}
\newcommand{\To}{\longrightarrow}
\newcommand{\BX}{\mathbf{B}(X)}
\newcommand{\A}{\mathcal{A}}
\newcommand{\SP}{\mathbb{P}^{S^1}}
\newcommand{\1}{{\textbf{\texttt{\tiny{1}}}}}
\newcommand{\2}{\textbf{\texttt{\tiny{2}}}}
\newcommand{\jj}{\textbf{\texttt{\tiny{j}}}}
\newcommand{\3}{\textbf{\texttt{\tiny{3}}}}
\newcommand{\ii}{\textbf{\texttt{\tiny{i}}}}
\newcommand{\kkk}{\textbf{\texttt{\tiny{k}}}}
\newcommand{\rrr}{\textbf{\texttt{\tiny{r}}}}
\newcommand{\ttt}{\textbf{\texttt{\tiny{t}}}}
\newcommand{\s}{\textbf{\texttt{\tiny{s}}}}
\newcommand{\ddd}{\textbf{\texttt{\tiny{d}}}}
\newcommand{\smallhollowcolon}{%
  \vcenter{\offinterlineskip
    \ialign{##\cr
      \tiny$\circ$\cr
      \noalign{\kern0.4ex}
      \tiny$\circ$\cr
    }%
  }%
}

% 空心正常序乘积 :: 样式
\newcommand{\normord}[1]{%
  \mathopen{\smallhollowcolon}%
  #1%
  \mathclose{\smallhollowcolon}%
}
\newcommand{\dd}{\textbf{\texttt{\tiny{d}}}}
\newcommand{\ve}{\mathbb{v}}
\newcommand{\we}{\mathbb{w}}
\newcommand{\tr}{\mathrm{tr}}
\newcommand{\ue}{\mathbb{u}}
\newcommand{\ch}{\mathrm{ch}}
\newcommand{\permute}{\tau}
\DeclareMathOperator{\pf}{Pf}
\usepackage{simpler-wick}
\usepackage{simpler-wick}
\usetikzlibrary{decorations.markings}
% from https://tex.stackexchange.com/a/39282/121799
\tikzset{W->-/.style={decoration={
  markings,
  mark=at position 0.5*\pgfdecoratedpathlength+2pt with
  {\draw[-latex] (-2pt,0pt) -- (1pt,0pt);}},postaction={decorate}},
  W-<-/.style={decoration={
  markings,
  mark=at position 0.5*\pgfdecoratedpathlength with
  {\draw[latex-] (-2pt,0pt) -- (1pt,0pt);}},postaction={decorate}}
  }
\pgfkeys{
  /simplerwick/.cd,
  arrows/.store in=\LstWickArrows,
  arrows/.initial={-,-,-,-,-,-,-,-,-}, % the # of contractions is bounded by 9
  arrows={-,-,-,-,-,-,-,-,-},
  positions/.store in=\LstWickPositions,
  positions={+1,+1,+1,+1,+1,+1,+1,+1,+1},
  positions/.initial={+1,+1,+1,+1,+1,+1,+1,+1,+1},
}

\makeatletter
\newcounter{Wick@up}
\newcounter{Wick@down}
\def\swick@end#1#2{
  \swick@setfalse@#1
  \tikzexternaldisable
  \begin{tikzpicture}[remember picture, baseline=(swick-close#1.base)]
    \node[use as bounding box, inner sep=0pt, outer sep=0pt] (swick-close#1) {$\displaystyle #2$};
  \end{tikzpicture}
  \tikz[remember picture, overlay]
{
\xdef\myW@style{\empty}
\foreach \W@X[count=\W@C] in \LstWickArrows
{\ifnum\W@C=#1
\xdef\myW@style{\W@X}
\fi}
\ifx\myW@style\empty
\PackageWarning{simpler-wick}{%
The list arrows has not enough entries!%
}{}
\xdef\myW@style{-}
\fi
\xdef\myW@pos{-77}
\foreach \W@X[count=\W@C] in \LstWickPositions
{\ifnum\W@C=#1
\xdef\myW@pos{\W@X}
\fi}
\ifnum\myW@pos=-77
\PackageWarning{simpler-wick}{%
The list positions has not enough entries!%
}{}
\xdef\myW@pos{+1}
\fi
\ifnum\myW@pos=-1
    \draw[\myW@style] ($(swick-open#1.south) + (0, -3pt)$) 
          -- ($(swick-open#1.base) + (0, -\swick@offset) + \theWick@down*(0, -\swick@sep)$) 
          -- ($(swick-close#1.base) + (0, -\swick@offset) + \theWick@down*(0, -\swick@sep)$) 
          -- ($(swick-close#1.south) + (0, -3pt)$);
\stepcounter{Wick@down}
\else
\stepcounter{Wick@up}
    \draw[\myW@style] ($(swick-open#1.north) + (0, 3pt)$) 
          -- ($(swick-open#1.base) + (0, \swick@offset) + \theWick@up*(0, \swick@sep)$) 
          -- ($(swick-close#1.base) + (0, \swick@offset) + \theWick@up*(0, \swick@sep)$) 
          -- ($(swick-close#1.north) + (0, 3pt)$);
\fi}
  \tikzexternalenable}
\def\wick@[#1]#2{\setcounter{Wick@up}{0}
\setcounter{Wick@down}{-1}
  \ifmmode
    \begingroup
    \pgfkeys{
        simplerwick,
        #1}
    % Define the variables and commands
    \swick@cond@reset
    \swick@count=0
    \def\swick@max{0}
    \def\c{\swick@smart}
    % Here is the text
    #2
    % Add a vbox equal to max height
    \dimen0=\swick@sep
    \multiply\dimen0 by \swick@max
    \advance\dimen0 by \swick@offset
    \vbox to \dimen0{}
    % Check that every has been closed
    \swick@cond@any{
      \PackageWarning{simpler-wick}{%
        I have reached the end of \protect\wick\space with some unclosed
        contractions%
      }{}
    }{}
    \endgroup
  \else
    \PackageWarning{simpler-wick}{%
      \protect\wich\space has been called outside a math environment, this will
      be ignore%
    }
  \fi
}

\makeatother
%\newcommand{\brian}[1]{(\textcolor{violet}{Note: #1})}

%\newcommand{\gui}[1]{(\textcolor{blue}{ZG: #1})}

\begin{document}
\maketitle

\abstract{
  We classify central extensions of the dg Lie algebra of derived global sections of the tangent sheaf on the
punctured, formal $d$-disk for $d > 2$. 
The case of $d=2$ has been handled in \cite{GWvir2d}.
We also prove a local, universal form of the Grothendieck--Riemann--Roch theorem for families of $d$-dimensional complex
varieties.}
% ----------------------------------------------------------------

%\tableofcontents

\section{Introduction}

Let $\mathring{D}^d$ denote the punctured $d$-dimensional formal disk.
For $d > 1$ this is not a formal affine scheme.
The main object of study in this paper is the derived global sections of the tangent sheaf 
\begin{equation}\label{}
  \lie{witt}_d \simeq \R \Gamma(\mathring{D}^d , \sT) 
\end{equation}
on the punctured $d$-disk.
We refer to this dg Lie algebra as the $d$-dimensional Witt algebra.
We write equivalence, rather than definition, since we use $\lie{witt}_d$ to refer to a concrete dg model for the derived global
sections which was introduced in \cite{FHK,GWWchiral} motivated by the theory of Jouanolou torsors.
See \cite{GWvir2d} for more details on the definition of $\lie{witt}_d$; we review its features in section
\ref{s:background}.

The first main result of this paper is the classification of central extensions of this dg Lie algebra.
We will show, see theorem \ref{thm:local}, that the second Lie algebra (hyper) cohomology of $\lie{witt}_d$ is $p(d+1)-1$
dimensional where $p(k)$ is the number of partitions of the integer $k$.
In fact, we use a universal Chern--Weil homomorphism \cite{GWvir2d}:
\begin{equation}\label{}
  \mathsf{cw}_d \colon \C[\op{ch}_1,\op{ch}_2, \ldots, \op{ch}_{d}] \to \HH^2_{Lie} (\lie{witt}_d) 
\end{equation}
which we will show is an isomorphism when restricted to polynomials in the $\op{ch}_i$ which are of total degree $d+1$.
Here, $\op{ch}_i$ is treated as a degree $i$ polynomial variable, and plays the role of the universal $i$th Chern character.
The proof proceeds in two steps: (1) first we will show in section \ref{s:chern} that $\mathsf{cw}_d$ is injective when restricted to degree
$d+1$ polynomials, (2) we next use a spectral sequence induced from the "diagonal filtration" in section \ref{s:gf} to obtain
an upper bound for the second cohomology.

Our next result is a local, universal Grothendieck--Riemann--Roch theorem.
To state it, we start with a tensor $\lie{witt}_d$-module $\cV$.
Roughly, this is a dg module obtained as the global sections of some natural vector bundle on the formal punctured
disk.
Now, such a tensor module defines a homomorphism 
\begin{equation}\label{}
  \mathsf{L}_V \colon \lie{witt}_d \to \op{End}(\cV) . 
\end{equation}
On the other hand, $\op{End}(\cV)$ is a dg Tate vector space. 
It follows \cite{FHK} that its second Lie algebra cohomology is one-dimensional generated by the preferred class 
\begin{equation}\label{}
  [z^\1 \wedge \cdots \wedge z^{\dd} \wedge P] \in \HH^2_{Lie}(\op{End}(\cV)) 
\end{equation}
where $[P] \in H^{d-1}(\mathring{D}^d, \sO)$ is the $d$-dimensional residue class.
We will prove in section \ref{s:GRR} that
\begin{equation}\label{}
  \mathsf{L}_V^* [z^\1 \wedge \cdots \wedge z^{\dd} \wedge P] = \op{Td} \op{Ch}(\cV) \in \HH^2_{Lie}(\lie{witt}_d) .
\end{equation}

\paragraph
Most of the essential definitions and constructions used in this paper have appeared in \cite{GWvir2d} where we proved versions of the main
results of this paper in the case $d=2$.
We recall the neccesary notations in \ref{s:background} and \ref{s:chern}, but we refer the reader to \textit{loc.
cit.} for more details.

\paragraph
In the main text we also study an extension of dg Lie algebra $\lie{witt}_d$ by the dg Lie algebra of $\lie{gl}
_r$-valued currents on the punctured $d$-disk.
This is $\lie{gl}_r \otimes \cJ_d$ where $\cJ_d \simeq \R \Gamma(\mathring{D}^d, \sO)$ is the Jouanolou model for the
structure sheaf.
On its own, central extensions of this current dg Lie algebra were introduced as higher-dimensional
Kac--Moody algebras in \cite{FHK}.
Our work extends the classification of central exstensions for the semi-direct procut $\lie{witt}_d \ltimes \lie{gl}
_r (\cJ_d)$.
We also formulate and prove an extension of the Grothendieck--Riemann--Roch theorem to this setting.

\subsection*{Acknowledgements}
We thank Kevin Costello, Mikhail Kapranov, and Matt Szczesny for valuable discussions, particularly during the early stages of this work.

\subfile{background}
\subfile{chern}
\subfile{gf}
\subfile{GRR}

\newpage 

\printbibliography
%\bibliographystyle{amsplain}
%\bibliography{RRChiral}
\end{document}




